% META: {"id": "TEXP-CTP-CO-010", "chapitre": "TEXP-COMPLEXES-TRIGO-POLYNOMES", "type_objet": "corrige", "exercice_ref": "TEXP-CTP-EX-010", "capacites_codes": ["C7"], "sources_inspiration": [], "status": "generated"}
% BEGIN-VERIFY
% from sympy import I, exp, pi, simplify, expand, symbols, Abs, arg, N, Poly
% z_sym = symbols('z_sym')
% racines = [exp(2 * I * k * pi / 5) for k in range(5)]
% for w in racines:
%     assert simplify(w ** 5 - 1) == 0
%     assert simplify(Abs(w) - 1) == 0
% # La somme est nulle par Vieta : dans z^5 - 1, le coefficient de z^4
% # est nul, et c'est l'oppose de la somme des racines. On le lit sur le
% # polynome, et on le confirme numeriquement a haute precision.
% assert Poly(z_sym ** 5 - 1, z_sym).all_coeffs()[1] == 0
% assert abs(complex(N(sum(racines), 50))) < 1e-40
% z = symbols('z')
% assert expand((z - 1) * (z ** 4 + z ** 3 + z ** 2 + z + 1)) == expand(z ** 5 - 1)
% # L'ecart angulaire constant se lit sur le QUOTIENT de deux racines
% # consecutives, qui vaut exactement exp(2 i pi / 5) : c'est exact, la ou
% # une soustraction d'arguments depend de la determination choisie.
% for k in range(4):
%     assert simplify(racines[k + 1] / racines[k] - exp(2 * I * pi / 5)) == 0
% END-VERIFY
\begin{corrige}{TEXP-CTP-EX-010}
\textbf{1.} Ce sont les $\omega_k=\mathrm{e}^{2\mathrm{i}k\pi/5}$ pour
$k=0,1,2,3,4$.

\textbf{2.} La somme vaut $0$. En effet
$z^5-1=(z-1)(z^4+z^3+z^2+z+1)$, et les $\omega_k$ pour $k\neq0$ annulent le
second facteur ; la somme des cinq racines est l'opposé du coefficient de
$z^4$ dans $z^5-1$, c'est-à-dire $0$.

\textbf{3.} Ces cinq points sont sur le cercle unité, régulièrement espacés
d'un angle $\tfrac{2\pi}{5}$ : ils forment un pentagone régulier inscrit dans
le cercle de centre $O$ et de rayon $1$.
\end{corrige}
